\chapter{वृहत्-विहित समुदाय}\label{ch:b3:grand-canonical}

आर्द्र हवा में पानी की एक बूँद न बढ़ती है न घटती; ऑक्सीजन फेफड़ों में
हीमोग्लोबिन से बँधती है और माँसपेशियों में छूट जाती है; और कोई बैटरी
इलेक्ट्रॉनों को फ़ोन से होकर इसलिए चलाती है कि वे एक इलेक्ट्रोड पर
दूसरे से ``अधिक महँगे'' हैं। हर बार जो मुद्रा तुलती है वह ताप नहीं,
बल्कि \emph{\omterm{thm:b3:grand-canonical:distribution}{रासायनिक विभव}} $\mu$ है: एक और कण की मुक्त-ऊर्जा लागत। यह
अध्याय दूसरा भंडार-वाल्व खोलता है --- ऊर्जा के साथ-साथ कण भी ---
और \omterm{thm:b3:grand-canonical:distribution}{वृहत्-विहित समुदाय} गढ़ता है, जिसका पुरस्कार है
चकित करने वाली दक्षता: किसी एक क्वांटम स्तर का अधिभोग, चाहे
फ़र्मिऑन हो या बोसॉन, दो पंक्तियों में निकल आता है, और उसी के साथ
अगले अध्याय के मूल सूत्र भी। इसके उपयोग गैस-मुखौटों और उत्प्रेरक
परिवर्तकों के अधिशोषण-समतापियों से लेकर
\omterm{prop:b3:grand-canonical:saha}{साहा समीकरण} तक फैले हैं, जिससे 1925 में हार्वर्ड की एक शोध-छात्रा ने
तारों के ताप पढ़े और, सारी अपेक्षाओं के विरुद्ध, यह खोजा कि
ब्रह्मांड हाइड्रोजन का बना है।

\section{बातचीत में कण भी शामिल}

\begin{theorem}[वृहत्-विहित वितरण]\label{thm:b3:grand-canonical:distribution}
\index{वृहत्-विहित समुदाय}\index{रासायनिक विभव}
कोई तंत्र, जो ताप $T$ और \omterm{thm:b3:grand-canonical:distribution}{रासायनिक विभव} $\mu$ वाले किसी भंडार से ऊर्जा
\emph{और कण} दोनों का विनिमय करता है, अपनी अवस्था $s$
(ऊर्जा $E_s$, कण संख्या $N_s$) में इस प्रायिकता से रहता है
\[
\mathcal P_s = \frac{\eu^{-\beta(E_s - \mu N_s)}}{\Xi} , \qquad
\Xi = \sum_s\eu^{-\beta(E_s - \mu N_s)} ,
\]
अर्थात् \emph{वृहत् \omterm{thm:b3:canonical-ensemble:boltzmann}{विभाजन फलन}}। $\Xi$ से: $\langle N\rangle
= k_{\text{B}}T\,\partial\ln\Xi/\partial\mu$, और \emph{वृहत्
विभव} $\Phi = -k_{\text{B}}T\ln\Xi$ $\Phi = -PV$ मानता है।
\end{theorem}

\begin{proof}[आंशिक उपपत्ति]
\cref{thm:b3:canonical-ensemble:boltzmann} की तरह, भंडार की
एन्ट्रॉपी को दोनों निकासियों में प्रसारित कीजिए:
$S_{\text{भं}}(E - E_s, N - N_s) \approx S_{\text{भं}} - E_s/T +
\mu N_s/T$, जहाँ $\partial S/\partial N = -\mu/T$
(\cref{thm:b3:microcanonical-ensemble:temperature})। औसत
अवकलन से निकलते हैं; और $\Phi = -PV$ मान लिया गया है (वह
विस्तीर्णता से निकलता है)।
\end{proof}

\begin{definition}[$\mu$ का अर्थ]\label{def:b3:grand-canonical:mu}
\index{रासायनिक विभव!अर्थ}
$\mu$ स्थिर $T$ और $V$ पर एक कण जोड़ने की मुक्त-ऊर्जा लागत है। कण
स्वतःस्फूर्त रूप से उच्च $\mu$ से निम्न
$\mu$ की ओर बहते हैं --- विसरण, परासरण, वाष्पन, रासायनिक
अभिक्रियाएँ और विद्युत् धाराएँ, सब ढीली पड़ती $\mu$-प्रवणताएँ हैं।
किन्हीं दो प्रावस्थाओं या स्थानों के बीच संतुलन पर $\mu$ दोनों ओर
बराबर होता है --- वही तीसरी समता ($T$ और $P$ के साथ), जो सह-अस्तित्व
पर शासन करती है। किसी
चिरसम्मत आदर्श गैस के लिए,
\[
\mu = k_{\text{B}}T\,\ln\!\big(n\lambda_T^3\big) < 0
\quad\text{(विरल गैस)} :
\]
किसी विरल गैस में एक कण जोड़ना \omterm{prop:b3:canonical-ensemble:machine}{मुक्त ऊर्जा} को \emph{घटा} देता है,
क्योंकि एन्ट्रॉपी का लाभ किसी भी ऊर्जा-लागत पर भारी पड़ता है --- वही
चिह्न, जो सूखी हवा में वाष्पन चलाता है। और कोई बैटरी एक $\mu$-मशीन
है: उसकी वोल्टता उसके इलेक्ट्रोडों के बीच प्रति इलेक्ट्रॉन
रासायनिक-विभव का पतन है, $U = \Delta\mu/e$ --- वोल्ट इलेक्ट्रॉन-रसायन
की विनिमय-दर हैं, और इसीलिए सेल एक-अंकीय वोल्ट देते हैं।
\end{definition}

\section{एक बार में एक स्तर: मूल तरकीब}

\begin{theorem}[किसी एक क्वांटम स्तर का अधिभोग]\label{thm:b3:grand-canonical:occupation}
\index{फ़र्मी--डिराक वितरण}\index{बोस--आइंस्टाइन वितरण}
``तंत्र'' के रूप में ऊर्जा $\epsilon$ का कोई एक एक-कण स्तर लीजिए, जो
बाकी सबकी गैस से वृहत् संपर्क में हो।
\emph{फ़र्मिऑन} (अधिभोग $0$ या $1$):
\[
\Xi = 1 + \eu^{-\beta(\epsilon - \mu)}
\;\Longrightarrow\;
\bar n_{\text{FD}} = \frac{1}{\eu^{\beta(\epsilon - \mu)} + 1} .
\]
\emph{बोसॉन} (अधिभोग $0, 1, 2, \dots$):
\[
\Xi = \frac{1}{1 - \eu^{-\beta(\epsilon - \mu)}}
\;\Longrightarrow\;
\bar n_{\text{BE}} = \frac{1}{\eu^{\beta(\epsilon - \mu)} - 1} ,
\]
जिसके लिए $\mu < \epsilon_{\min}$ चाहिए। अधिभोग विरल होने पर दोनों
बोल्ट्ज़मान गुणक $\eu^{-\beta(\epsilon - \mu)}$ पर सिमट आते हैं।
ये दो सूत्र, जो दो-दो पंक्तियों में मिले, क्वांटम सांख्यिकी का पूरा
निवेश हैं: अगला अध्याय इन्हीं की फ़सल है।
\end{theorem}

\begin{proof}
फ़र्मिऑन: दो पद। बोसॉन: एक गुणोत्तर श्रेणी, जो केवल
$\mu < \epsilon$ के लिए अभिसरित होती है। दोनों ही स्थितियों में $\bar n = k_{\text{B}}T\,
\partial\ln\Xi/\partial\mu$ कथित रूप दे देता है।
\end{proof}

\begin{omfigure}
\begin{tikzpicture}
  \begin{axis}[omaxis, axis x line=bottom, axis y line=left,
    width=11.2cm, height=5.6cm, xmin=-6, xmax=6, ymin=0, ymax=2.3,
    xlabel={$(\epsilon - \mu)/k_{\text{B}}T$}, ylabel={$\bar n$},
    xlabel style={at={(axis description cs:0.5,-0.1)}, anchor=north},
    xtick={-4,-2,0,2,4}, ytick={0.5,1,2}, clip=false,
    legend style={font=\scriptsize, at={(0.97,0.93)}, anchor=north east,
      draw=none, fill=none}]
    \addplot[omDef, very thick, domain=-6:6, samples=200]
      {1/(exp(x)+1)};
    \addlegendentry{फ़र्मी--डिराक}
    \addplot[omThm, very thick, domain=0.45:6, samples=200]
      {1/(exp(x)-1)};
    \addlegendentry{बोस--आइंस्टाइन}
    \addplot[omExo, dashed, thick, domain=-0.8:6, samples=100]
      {exp(-x)};
    \addlegendentry{बोल्ट्ज़मान}
    \draw[gray, dashed] (axis cs:0,0) -- (axis cs:0,2.3);
  \end{axis}
\end{tikzpicture}

{\small तीनों अधिभोग-नियम, $(\epsilon -
\mu)/k_{\text{B}}T$ के सामने: फ़र्मिऑन प्रति अवस्था एक पर संतृप्त हो
जाते हैं ($\epsilon = \mu$ पर एक चिकनी सीढ़ी), बोसॉन
$\epsilon \to \mu$ पर बिना सीमा ढेर होते जाते हैं, और दाईं ओर दोनों
विरल बोल्ट्ज़मान पूँछ में मिल जाते हैं --- चिरसम्मत प्रांत वहीं है
जहाँ किसी अवस्था में शायद ही कभी जाया जाता हो।}
\end{omfigure}

\section{अधिशोषण: किराए पर स्थल}

\begin{proposition}[लांगम्यूर समतापी]\label{prop:b3:grand-canonical:langmuir}
\index{लांगम्यूर समतापी}\index{अधिशोषण}
कोई सतह स्वतंत्र स्थल देती है, जिनमें से हर एक या तो खाली है या
बंधन ऊर्जा $-\epsilon_0$ के साथ एक गैस-अणु रखता है; और ऊपर की गैस,
दाब $P$ पर, भंडार है। हर स्थल एक द्वि-अवस्था वृहत्
तंत्र है, और उसका अधिभोग है
\[
\theta = \frac{1}{\eu^{\beta(-\epsilon_0 - \mu)} + 1}
 = \frac{P}{P + P_0(T)} , \qquad
P_0(T) = \frac{k_{\text{B}}T}{\lambda_T^3}\,
\eu^{-\epsilon_0/k_{\text{B}}T} ,
\]
आदर्श-गैस वाले $\mu$ का उपयोग करके। आच्छादन कम दाब पर रैखिक रूप से
बढ़ता है और एक एकपरत पर संतृप्त हो जाता है --- अर्थात् \emph{\omterm{prop:b3:grand-canonical:langmuir}{लांगम्यूर
समतापी}}, जो उत्प्रेरक परिवर्तकों, सक्रियित
कोयले, गैस-मुखौटों और रक्त की ऑक्सीजन-बंधन का कार्यकारी वक्र है
(\cref{exo:b3:grand-canonical:6})। संक्रमण-दाब $P_0$
बंधन ऊर्जा के साथ \emph{चरघातांकी} रूप से गिरता है और ताप के साथ
तेज़ी से चढ़ता है: किसी सतह को गरम कीजिए और उसके मेहमान झड़ जाते हैं
--- इसीलिए निर्वात में काँच के पात्र सेंकना ही उन्हें सचमुच साफ़ करना है।
\end{proposition}

\begin{proof}
स्थल का अधिभोग
\cref{thm:b3:grand-canonical:occupation} है (फ़र्मिऑनी रूप: कोई स्थल
अधिक से अधिक एक रखता है), जहाँ $\epsilon = -\epsilon_0$; और
$\eu^{\beta\mu} = n\lambda_T^3 = P\lambda_T^3/k_{\text{B}}T$ रखिए।
\end{proof}

\begin{omfigure}
\begin{tikzpicture}
  \begin{axis}[omaxis, axis x line=bottom, axis y line=left,
    width=10.8cm, height=5.4cm, xmin=0, xmax=5, ymin=0, ymax=1.12,
    xlabel={निम्नतर ताप पर $P/P_0$}, ylabel={आच्छादन $\theta$},
    xlabel style={at={(axis description cs:0.5,-0.1)}, anchor=north},
    xtick={1,2,3,4}, ytick={0.5,1}, clip=false,
    legend style={font=\scriptsize, at={(0.97,0.42)}, anchor=north east,
      draw=none, fill=none}]
    \addplot[omDef, very thick, domain=0:5, samples=120] {x/(x+1)};
    \addlegendentry{ठंडी सतह}
    \addplot[omThm, very thick, domain=0:5, samples=120] {x/(x+6)};
    \addlegendentry{गरम सतह ($P_0$ बड़ा)}
    \draw[gray, dashed] (axis cs:0,1) -- (axis cs:5,1);
    \node[font=\scriptsize, anchor=south east, gray] at (axis cs:4.9,1.0)
      {एक एकपरत};
  \end{axis}
\end{tikzpicture}

{\small \omterm{prop:b3:grand-canonical:langmuir}{लांगम्यूर समतापी}: दाब के सामने आच्छादन, जो
एक एकपरत पर संतृप्त होता है। सतह गरम कीजिए और वही आच्छादन कहीं
अधिक दाब माँगता है --- अधिशोषण एक $\mu$-सौदेबाज़ी है, जिसे गरम सतह
हारती रहती है।}
\end{omfigure}

\section{आयनन संतुलन: साहा समीकरण}

\begin{proposition}[साहा का समीकरण]\label{prop:b3:grand-canonical:saha}
\index{साहा समीकरण}
किसी तप्त गैस में अभिक्रिया $\text{H} \rightleftharpoons \text{p} +
\text{e}$ तब संतुलित होती है जब $\mu_{\text{H}} = \mu_{\text{p}} +
\mu_{\text{e}}$। हर $\mu$ को आदर्श-गैस रूप में लिखने पर (और परमाणु
के खाते में उसका बंधन $-E_{\text{I}}$ जमा करके) मिलता है
\[
\frac{n_{\text{e}}\,n_{\text{p}}}{n_{\text{H}}}
 = \frac{1}{\lambda_{T,\text{e}}^3}\,
\eu^{-E_{\text{I}}/k_{\text{B}}T}
\qquad\Big(\lambda_{T,\text{e}} =
h/\sqrt{2\pi m_{\text{e}}k_{\text{B}}T}\Big)
\]
(कोटि-एक के चक्रण-अपह्रासिता गुणकों को छोड़कर)। फिर वही दो
प्रतिस्पर्धी मुद्राएँ: बोल्ट्ज़मान गुणक आयनन पर
$E_{\text{I}} = \qty{13.6}{eV}$ का कर लगाता है, जबकि अगुणक --- अर्थात्
\emph{मुक्त हुए इलेक्ट्रॉन की एन्ट्रॉपी}, अपनी $\lambda_T^{-3}$
जितनी सुलभ अवस्थाओं के साथ --- उसे अनुदान देता है। तारकीय घनत्वों पर
वह अनुदान इतना बड़ा है कि हाइड्रोजन \qty{10000}{K} के पास ही ठीक-ठाक
आयनित हो जाता है, जहाँ $k_{\text{B}}T$ केवल
\qty{0.9}{eV} है: संतुलन \omterm{prop:b3:canonical-ensemble:machine}{मुक्त ऊर्जा} पर लड़े जाते हैं, ऊर्जा पर नहीं
--- और तारों के वर्णक्रम ही उसका बहीखाता हैं
(\cref{pb:b3:grand-canonical:1})।
\end{proposition}

\begin{proof}[आंशिक उपपत्ति]
$\mu_{\text{H}} = \mu_{\text{p}} + \mu_{\text{e}}$ रखिए, जहाँ $\mu_i
= k_{\text{B}}T\ln(n_i\lambda_{T,i}^3) + \varepsilon_i$
($\varepsilon_{\text{H}} = -E_{\text{I}}$, शेष शून्य); और
चरघातांक लीजिए। प्रोटॉन तथा परमाणु की तापीय तरंगदैर्घ्य लगभग
कट जाती हैं (द्रव्यमान बराबर); और अपह्रासिता गुणक मान लिए गए हैं।
\end{proof}

\begin{omfigure}
\begin{tikzpicture}
  \begin{axis}[omaxis, axis x line=bottom, axis y line=left,
    width=10.8cm, height=5.4cm, xmin=3000, xmax=25000, ymin=0, ymax=1.1,
    xlabel={$T$ (K)}, ylabel={हाइड्रोजन का आयनित अंश},
    xlabel style={at={(axis description cs:0.5,-0.1)}, anchor=north},
    xtick={5000,10000,15000,20000}, ytick={0.5,1}, clip=false]
    \addplot[omDef, very thick, domain=3000:25000, samples=300]
      {1/(1+exp(-(x-9600)/1100))};
    \draw[gray, dashed] (axis cs:9600,0) -- (axis cs:9600,0.5);
    \node[font=\scriptsize, anchor=west] at (axis cs:10000,0.28)
      {\qty{10000}{K} के पास अर्ध-आयनित};
  \end{axis}
\end{tikzpicture}

{\small प्रकाशमंडलीय घनत्व पर हाइड्रोजन का साहा आयनन
(योजनाबद्ध): \qty{10000}{K} के पास एक तीखा संक्रमण, जो भोले
$E_{\text{I}}/k_{\text{B}} \approx \qty{160000}{K}$ से कहीं नीचे है
--- मुक्त हुए इलेक्ट्रॉन की एन्ट्रॉपी ही अधिकांश बिल चुकाती है।}
\end{omfigure}

\begin{method}[वृहत्-विहित कारीगरी]\label{met:b3:grand-canonical:method}
(1) जो कुछ भी किसी बड़े साथी से कण बदलता हो: साथी को कोई $\mu$ दीजिए
और अवस्थाओं को $\eu^{-\beta(E - \mu N)}$ से भारित कीजिए।
(2) स्वतंत्र स्थल या स्तर गुणनखंडित हो जाते हैं: \emph{एक} हल कीजिए,
और गुणा कर दीजिए। (3) किसी चिरसम्मत-गैस भंडार के लिए $\eu^{\beta\mu} =
n\lambda_T^3$ $\mu$ को मापनीय दाब में बदल देता है। (4) अभिक्रियाएँ
और प्रावस्था-परिवर्तन: हर जाति के लिए एक-एक $\mu$ का संतुलन बिठाइए,
बंधन ऊर्जाएँ सहित --- और अपेक्षा रखिए कि एन्ट्रॉपी-अगुणक संतुलनों को
भोले बोल्ट्ज़मान अनुमानों से बहुत दूर खिसका देंगे। (5) ठोसों में
इलेक्ट्रॉन: उनका $\mu$ ही फ़र्मी स्तर है, और $\mu$ के अंतर
वोल्टताएँ हैं --- यही \cref{ch:b3:electrons-in-solids} तक का पुल है।
\end{method}

\section{अभ्यास}

\begin{exercise}[$\star$]\label{exo:b3:grand-canonical:1}
(क) एक ही गैस के दो बक्से, एक ही $T$ पर, पर भिन्न घनत्वों वाले, जोड़
दिए जाते हैं: कण किस ओर बहते हैं,
$\mu$ की भाषा में? (ख) आदर्श-गैस वाले $\mu$ से दिखाइए कि बहाव
बराबर घनत्वों पर रुक जाता है। (ग) आर्द्र हवा के नीचे पानी रखा है:
संतुलन की शर्त $\mu$ के पदों में बताइए। (घ) हवा सूखी हो तो कपड़े
$100\,^\circ$C से नीचे भी क्यों सूख जाते हैं ($\mu$ की तुलना कीजिए,
तापों की नहीं)?
\end{exercise}

\begin{exercise}[$\star$]\label{exo:b3:grand-canonical:2}
हवा का \omterm{thm:b3:grand-canonical:distribution}{रासायनिक विभव}। (क) $\lambda_T$ निकालिए, N$_2$ के लिए
\qty{300}{K} पर। (ख) \qty{1}{bar} पर $n\lambda_T^3$ और उससे निकलता
$\mu$ eV में निकालिए। (ग) $\mu < 0$ कोई विरोधाभास क्यों नहीं है
(अणु जोड़ने पर $F = E - TS$ का कौन-सा पद हावी रहता है)? (घ)
किस घनत्व पर $\mu$ शून्य तक पहुँचेगा, और
$n\lambda_T^3 \sim 1$ किसकी आहट है
(\cref{exo:b3:grand-canonical:11})?
\end{exercise}

\begin{exercise}[$\star$]\label{exo:b3:grand-canonical:3}
फ़र्मी--डिराक अधिभोग को उसके दो-पदी $\Xi$ से
कदम-दर-कदम व्युत्पन्न कीजिए; सीमाएँ $\epsilon \ll \mu$, $\epsilon \gg
\mu$ और $\epsilon = \mu$ पर मान जाँचिए। यह फलन इस पुस्तक में
पहले कहाँ प्रकट हो चुका है
(\cref{exo:b3:microcanonical-ensemble:8})?
\end{exercise}

\begin{exercise}[$\star$]\label{exo:b3:grand-canonical:4}
वोल्ट रसायन हैं। (क) कोई लीथियम सेल \qty{3.0}{V} रखता है:
प्रति इलेक्ट्रॉन रासायनिक-विभव का अंतर eV में और प्रति मोल kJ में
व्यक्त कीजिए। (ख) रोज़मर्रा की सारी बैटरियाँ एक-अंकीय
वोल्ट की क्यों होती हैं (रासायनिक $\mu$-अंतरों का पैमाना क्या तय
करता है)? (ग) कोई सीसा--अम्ल कार-बैटरी छह सेल जोड़ती है: जोड़ने के
बदले कोई \qty{12}{eV} अभिक्रिया क्यों नहीं खोजी जाती? (घ) $\mu$ की
भाषा में, ``बैठ चुकी'' बैटरी क्या कर रही होती है?
\end{exercise}

\begin{exercise}[$\star\star$]\label{exo:b3:grand-canonical:5}
किसी स्तर पर बोसॉन। (क) $\Xi$ के लिए गुणोत्तर श्रेणी का योग कीजिए और
$\bar n_{\text{BE}}$ व्युत्पन्न कीजिए। (ख) $\mu$ को सबसे नीचे वाले
स्तर से नीचे क्यों रहना चाहिए --- अन्यथा क्या अपसरित होता है, और वह
अपसरण किसकी आहट देता है (\cref{ch:b3:quantum-statistics})? (ग)
$\bar n$ निकालिए, $(\epsilon - \mu)/k_{\text{B}}T = 0.1$, $1$,
$3$ के लिए, और बोल्ट्ज़मान से तुलिए। (घ) फ़ोटॉनों का $\mu = 0$ होता है:
भौतिक कारण दीजिए (उनकी संख्या संरक्षित नहीं है --- दीवारें उन्हें
स्वतंत्र रूप से सोखती और उत्सर्जित करती हैं)।
\end{exercise}

\begin{exercise}[$\star\star$]\label{exo:b3:grand-canonical:6}
लांगम्यूर काम पर। (क) समतापी से: कौन-सा दाब
$\theta = 0.5$ देता है? $0.9$? (ख) मायोग्लोबिन O$_2$ को एक ही स्थल
पर बाँधता है: संतृप्ति ऑक्सीजन के आंशिक दाब में लांगम्यूर का
अनुसरण करती है --- मायोग्लोबिन के लिए \emph{अतिपरवलयिक} (न कि
सिग्मॉइडी) वक्र क्यों दिखता है, जबकि चार-स्थल वाले हीमोग्लोबिन के लिए
कोई सहकारी S-वक्र? (ग) कार्बन मोनोऑक्साइड उन्हीं स्थलों से बँधती है,
जिसका $\epsilon_0$ $\approx\qty{0.13}{eV}$ बड़ा है: बराबर आंशिक दाबों
और \qty{310}{K} पर CO:O$_2$ अधिभोग-अनुपात निकालिए और CO विषाक्तता
समझाइए। (घ) समतापी की भाषा में, अति-दाब ऑक्सीजन चिकित्सा क्यों काम
करती है?
\end{exercise}

\begin{exercise}[$\star\star$]\label{exo:b3:grand-canonical:7}
संख्या के उतार-चढ़ाव। (क) दिखाइए कि $\operatorname{Var}N =
k_{\text{B}}T\,\partial\langle N\rangle/\partial\mu$। (ख) चिरसम्मत
आदर्श गैस के लिए निष्कर्ष निकालिए कि $\operatorname{Var}N = \langle
N\rangle$: अर्थात् पुआसों। (ग) इसे
\cref{exo:b3:microcanonical-ensemble:12} के घनत्व-उतार-चढ़ावों से
मिलाइए। (घ) किसी एक फ़र्मिऑनी स्तर के लिए
$\operatorname{Var}n = \bar n(1 - \bar
n)$ निकालिए: फ़र्मिऑन-संख्याएँ कहाँ सबसे शांत हैं, और क्यों?
\end{exercise}

\begin{exercise}[$\star\star$]\label{exo:b3:grand-canonical:8}
$\mu$ से परासरण। कोई झिल्ली पानी को जाने देती है, विलेय को नहीं
(सांद्रता $c$ प्रति आयतन, विरल)। (क) झिल्ली के दोनों ओर पानी का $\mu$
बराबर करके दिखाइए कि शुद्ध-पानी वाली ओर
$\Pi = ck_{\text{B}}T$ जितना अधिक दाब होना चाहिए (फ़ान 'ट हॉफ़) ---
यह मान लीजिए कि घुला विलेय प्रति जल-अणु पानी का $\mu$
$-k_{\text{B}}T\,c/n_{\text{w}}$ जितना घटा देता है। (ख)
$\Pi$ निकालिए, शरीर-क्रियात्मक लवण-जल ($c \approx \qty{3e26}{m^{-3}}$
कुल कण) के लिए। (ग) शुद्ध पानी में लाल रुधिर कोशिकाएँ: पानी किस ओर
बहता है, और परिणाम क्या? (घ) पेड़ों को दसियों मीटर तक रस चढ़ाने के
लिए न पंप चाहिए न चमत्कार --- क्यों (परासरणी और केशिकीय
$\mu$-प्रवणताएँ)?
\end{exercise}

\begin{exercise}[$\star\star$]\label{exo:b3:grand-canonical:9}
किसी क्षेत्र में संतुलन। गुरुत्व में किसी गैस-स्तंभ के संतुलन के लिए
\emph{कुल} \omterm{thm:b3:grand-canonical:distribution}{रासायनिक विभव} $\mu(h) =
k_{\text{B}}T\ln(n\lambda_T^3) + mgh$ का एकसमान होना ज़रूरी है। (क)
इससे वायुदाबमापी नियम एक पंक्ति में व्युत्पन्न कीजिए। (ख) कोणीय वेग
$\omega$ वाले किसी अपकेंद्रित्र के लिए व्यापक कीजिए (विभव
$-\tfrac12 m\omega^2r^2$) और त्रिज्य प्रोफ़ाइल पाइए। (ग)
यूरेनियम-संवर्धन के अपकेंद्रित्र $\omega r \sim
\qty{600}{m/s}$ पर घूमते हैं: प्रति चरण $^{238}$UF$_6$ के सामने
$^{235}$UF$_6$ के संवर्धन का संतुलन-अनुपात $\exp[\Delta m\,\omega^2
r^2/2k_{\text{B}}T]$ निकालिए, जहाँ $\Delta m = \qty{5e-27}{kg}$, $T =
\qty{320}{K}$। (घ) कई चरण शृंखला में क्यों लगाए जाते हैं?
\end{exercise}

\begin{exercise}[$\star\star\star$]\label{exo:b3:grand-canonical:10}
साहा, हाथ से। (क) \cref{prop:b3:grand-canonical:saha} की
$\mu$-संतुलन उपपत्ति कीजिए। (ख) आयनित अंश $x$ परिभाषित कीजिए, कुल
हाइड्रोजन घनत्व $n$ पर, और दिखाइए कि $x^2/(1 - x)
= \eu^{-E_{\text{I}}/k_{\text{B}}T}/n\lambda_{T,\text{e}}^3$। (ग)
सौर प्रकाशमंडल: $T = \qty{5800}{K}$, $n \approx
\qty{e23}{m^{-3}}$: दिखाइए कि हाइड्रोजन का आयनित अंश केवल
$\sim\num{e-4}$ है --- सूर्य की सतह उदासीन है। (घ) $T =
\qty{10000}{K}$ पर (वही $n$): फिर से निकालिए और तीखेपन पर टिप्पणी
कीजिए।
\end{exercise}

\begin{exercise}[$\star\star\star$]\label{exo:b3:grand-canonical:11}
जब गैसें क्वांटम हो जाती हैं। अपह्रासन $n\lambda_T^3 \sim
1$ पर शुरू होता है। (क) \qty{300}{K} की हवा के लिए: देहली से कितना
नीचे (\cref{exo:b3:grand-canonical:2})? (ख) ताँबे के चालन-इलेक्ट्रॉनों
($n = \qty{8.5e28}{m^{-3}}$) के लिए: वह ताप खोजिए जिसके नीचे
$n\lambda_T^3 > 1$ --- और निष्कर्ष निकालिए कि किसी धातु के
इलेक्ट्रॉन \emph{किसी भी} प्रयोगशाला-ताप पर अपह्रासित हैं।
(ग) $n = \qty{e20}{m^{-3}}$ पर रुबिडियम परमाणुओं के लिए: क्वांटम
देहली-ताप (\cref{ch:b3:quantum-statistics} का BEC पैमाना)। (घ)
तीनों को क्रम में रखिए और नियम बताइए: हलके कण और ऊँचे घनत्व पहले
क्वांटम होते हैं।
\end{exercise}

\begin{exercise}[$\star\star\star$]\label{exo:b3:grand-canonical:12}
ब्रह्मांड अपने पॉज़िट्रॉन खो देता है। तप्त आरंभिक ब्रह्मांड में
$\eu^+ + \eu^- \rightleftharpoons 2\gamma$ ने युग्मों को
$\mu_{e^+} + \mu_{e^-} = 0$ के साथ संतुलन में रखा (फ़ोटॉन: $\mu = 0$)।
(क) किसी सममित प्लाज़्मा के लिए इससे $\mu_{e^\pm} \approx 0$ मिलता है:
दिखाइए कि तब युग्म-घनत्व $n_\pm \approx
\lambda_T^{-3}\eu^{-m_{\text{e}}c^2/k_{\text{B}}T}$ है, जब एक बार
$k_{\text{B}}T \ll m_{\text{e}}c^2$ हो जाए। (ख) वह ताप निकालिए
जिस पर $m_{\text{e}}c^2 = k_{\text{B}}T$, और उससे संगत ब्रह्मांडीय
समय-पैमाना भी ($t \sim \qty{1}{s}$ मान लीजिए,
$\qty{e10}{K}$ पर)। (ग) दिखाइए कि थोड़ा और शीतलन (कुछ ही गुना)
युग्म-घनत्व को कई कोटियों से ढहा देता है:
पॉज़िट्रॉन नष्ट होकर लुप्त हो जाते हैं, और पीछे इलेक्ट्रॉनों का वही
नन्हा आदिम आधिक्य बचता है। (घ) मुक्त हुए फ़ोटॉन
विकिरण-कुंड को फिर गरम करते हैं: कौन-सा अवशेष विकिरण उसकी रसीद ढोता
है (\cref{ch:b3:astrophysics})?
\end{exercise}

\section{समस्या: तारों के ताप पढ़ना}

\begin{problem}[{सप्ताहांत समस्या --- साहा, पेन, और ब्रह्मांड
किससे बना है}]\label{pb:b3:grand-canonical:1}
तारों के वर्णक्रम 1890 के दशक में अपनी अवशोषण-रेखाओं की प्रबलता के
अनुसार वर्गों O, B, A, F, G, K, M में सूचीबद्ध किए गए थे ---
जिनमें हाइड्रोजन की बामर रेखाएँ वर्ग A में सबसे प्रबल थीं और उससे
अधिक गरम तथा अधिक ठंडे, दोनों तरह के तारों में रहस्यमय ढंग से
कमज़ोर। 1925 में सेसीलिया पेन ने बिलकुल नया \omterm{prop:b3:grand-canonical:saha}{साहा समीकरण} लगाया और
दिखाया कि वह पूरा चिड़ियाघर भिन्न तापों पर एक ही पदार्थ है --- और यह
भी कि तारे अत्यधिक रूप से हाइड्रोजन हैं, ऐसा निष्कर्ष जो इतना
विधर्मी था कि उन्हें अपने शोध-प्रबंध में उसे नरम करना पड़ा। यह समस्या
उनका तर्क फिर से चलती है।
आँकड़े: बामर अवशोषण $n = 2$ से शुरू होता है, जो मूल अवस्था से $E_2 - E_1 =
\qty{10.2}{eV}$ ऊपर है; $E_{\text{I}} =
\qty{13.6}{eV}$; प्रकाशमंडलीय इलेक्ट्रॉन घनत्व $n_{\text{e}}
\sim \qty{e20}{m^{-3}}$।

\medskip\noindent\textbf{भाग I --- दो प्रतिस्पर्धी गुणक।}
\begin{enumerate}
  \item बामर अवशोषण को उदासीन हाइड्रोजन चाहिए, वह भी $n =
        2$ \emph{में}: $n = 2$ की जनसंख्या के लिए बोल्ट्ज़मान गुणक
        लिखिए (अपह्रासिताएँ: $g_2/g_1 = 4$)।
  \item उसका मान $5000$, $10000$ और \qty{20000}{K} पर निकालिए: $T$
        बढ़ने पर यह गुणक रेखा-प्रबलता को किस ओर धकेलता है?
  \item अब साहा गुणक: $T$ बढ़ने पर \emph{उदासीन} अंश का क्या होता
        है, और इस तरह बामर रेखाओं का?
  \item गुणात्मक रूप से समझाइए कि गुणनफल --- उत्तेजित \emph{और}
        फिर भी उदासीन --- किसी मध्यवर्ती ताप पर शिखर क्यों बनाना
        ही चाहिए।
  \item \omterm{prop:b3:grand-canonical:saha}{साहा समीकरण} को $n_{\text{e}} =
        \qty{e20}{m^{-3}}$ के साथ काम में लेकर हाइड्रोजन के
        अर्ध-आयनन का ताप आकलिए।
  \item शिखर जोड़िए: बामर रेखाएँ किस ताप के पास सबसे प्रबल होनी
        चाहिए, और वह कौन-सा वर्णक्रमीय वर्ग
        (A तारे, $\sim\qty{10000}{K}$) चुनता है?
\end{enumerate}

\medskip\noindent\textbf{भाग II --- वर्णक्रमीय क्रम का अर्थ
खुला।}
\begin{enumerate}[resume]
  \item M तारे (\qty{3000}{K}) प्रबल आणविक बैंड (TiO) और कमज़ोर
        हाइड्रोजन दिखाते हैं: दोनों को इन्हीं दो गुणकों से समझाइए।
  \item O तारे (\qty{40000}{K}) हीलियम रेखाएँ और फिर से कमज़ोर
        हाइड्रोजन दिखाते हैं: समझाइए (हाइड्रोजन का क्या हुआ, और
        हीलियम को, जिसका $E_{\text{I}} = \qty{24.6}{eV}$ है, अब
        बारी क्यों मिली?)।
  \item सूर्य (G, \qty{5800}{K}) \emph{आयनित कैल्सियम}
        ($E_{\text{I,Ca}} = \qty{6.1}{eV}$) की प्रबल रेखाएँ दिखाता
        है, फिर भी कमज़ोर हाइड्रोजन: साहा-तर्क से दिखाइए कि
        कैल्सियम पहले ही आयनित क्यों है, जबकि हाइड्रोजन अब भी
        उदासीन और अधिकतर $n = 1$ में है।
  \item पेन से पहले के खगोलविद रेखा-प्रबलता को प्रचुरता पढ़ते थे:
        सूर्य की भीमकाय कैल्सियम रेखाओं और दुर्बल बामर रेखाओं ने
        उन्हें उसकी संरचना के बारे में क्या निष्कर्ष दिलाया?
  \item पेन की सूझ: हर रेखा की प्रबलता को उसके बोल्ट्ज़मान--साहा
        ``दृश्यता गुणक'' के लिए संशोधित कीजिए। जब उन्होंने ऐसा
        किया, तो सच्ची प्रचुरताओं का कौन-सा क्रम उभरा?
  \item उनका परिणाम --- हाइड्रोजन कैल्सियम से दस लाख गुना अधिक
        प्रचुर, और तारे अधिकतर हाइड्रोजन तथा हीलियम --- उसका विरोध
        क्यों हुआ (पृथ्वी-केंद्रित रसायन ने क्या सुझाया था), और
        उन्हें सही कब माना गया?
\end{enumerate}

\medskip\noindent\textbf{भाग III --- शिखर पर संख्याएँ।}
\begin{enumerate}[resume]
  \item ठीक \qty{9600}{K} पर $n = 2$ का बोल्ट्ज़मान गुणक
        निकालिए।
  \item मद 5 के साहा आकलन के साथ, $9600$, $15000$ और \qty{25000}{K}
        पर उदासीन अंश को मोटे तौर पर $0.5$, $\num{e-3}$,
        $\num{e-5}$ मान लीजिए: अपने तीनों तापों पर गुणनफल
        (उत्तेजित अंश $\times$ उदासीन अंश) बनाइए और उसका अधिकतम
        खोजिए।
  \item बामर शिखर गरम ओर \emph{तीखा} (कौन-सा चरघातांकी जीतता है)
        और ठंडी ओर \emph{क्रमिक} क्यों है?
  \item किसी श्वेत-वामन प्रकाशमंडल का $n_{\text{e}}$ हज़ार गुना
        बड़ा होता है: अर्ध-आयनन का ताप किस ओर खिसकता है, और क्यों
        ($\mu$ की भाषा में ल शातलिए)?
  \item अतः एक ही ताप के, पर भिन्न दाबों वाले दो तारे भिन्न
        रेखा-प्रबलताएँ दिखाते हैं: इससे वर्णक्रममितिज्ञ तारे का
        कौन-सा गुण (दानव बनाम वामन --- सतही गुरुत्व) पढ़ लेते हैं?
  \item आधुनिक सर्वेक्षण रेखा-अनुपातों से तारकीय ताप $\pm1\%$ तक
        बाँध देते हैं: शृंखला बताइए --- अभिगृहीत $\to$ बोल्ट्ज़मान
        $\to$ साहा $\to$ तापमापी।
\end{enumerate}

\medskip\noindent\textbf{भाग IV --- इससे क्या तय हुआ।}
\begin{enumerate}[resume]
  \item OBAFGKM क्रम आनुभविक रूप से पुनर्व्यवस्थित वर्णमालाई
        अव्यवस्था था: साहा के बाद कौन-सा एक भौतिक चर उसे क्रम में
        बिठाता है?
  \item बताइए कि \emph{वही} तत्त्व किसी एक ताप पर वर्णक्रम पर छा
        सकता है और किसी दूसरे पर लुप्त हो सकता है, बिना प्रचुरता
        में किसी परिवर्तन के --- वही मूल सीख, जो पेन ने
        खगोल-विज्ञान को दी।
  \item हाइड्रोजन सारे बैरिऑनी द्रव्यमान का तीन-चौथाई: इस पुस्तक के
        दो और स्तंभ बताइए जो उसी प्रचुरता पर टिके हैं (तारकीय
        संलयन, \cref{ch:b3:astrophysics}; 21 cm का
        आकाश, \cref{ch:b3:spin-two-level})।
  \item ओटो श्ट्रूफ़े ने पेन के शोध-प्रबंध को ``खगोल-विज्ञान में अब
        तक लिखा गया सबसे प्रतिभाशाली PhD शोध-प्रबंध'' कहा: इस
        समस्या से उस प्रशंसा को एक वाक्य में उचित ठहराइए।
  \item वही साहा भौतिकी, ब्रह्मांडीय पैमाने पर उलटी चलाई जाए, तो
        वह ताप बताती है जिस पर पूरे ब्रह्मांड का हाइड्रोजन पहली बार
        जुड़ा (कहीं कम ब्रह्मांडीय घनत्व पर $\sim\qty{3000}{K}$):
        उस घटना और उसके अवशेष का नाम लीजिए
        (\cref{ch:b3:astrophysics})।
  \item आयनन संतुलन प्रतिदीप्त लैंप और प्लाज़्मा-प्रसंस्करण भी
        चलाते हैं: एक वाक्य में बताइए कि किसी ऐसी नली में, जिसका
        काँच ठंडा बना रहता है, \qty{10000}{K}-सरीखी आयनित अवस्था
        क्यों हो सकती है (कौन-से दो उप-तंत्र संतुलन में नहीं आते)?
  \item नामित परिणाम का सार दीजिए: एक ही संतुलन-शर्त,
        $\mu_{\text{H}} = \mu_{\text{p}} + \mu_{\text{e}}$,
        वर्णक्रमीय वर्णमाला को $3000$ से \qty{40000}{K} तक चलने
        वाला तापमापी बना देती है, बामर शिखर को A तारों पर रखती है,
        और हाइड्रोजन से बना एक ब्रह्मांड उजागर करती है।
\end{enumerate}
\end{problem}
